\chapter{Introduction}

Agricultural fields require repeated observation over large area. UAV swarms provide flexible aerial imaging solutions for crop monitoring, field inspection, and early detection of abnormal conditions of large agricultural fields. However UAVs have constraints of energy, computation and connectivity, real-time CNN inference remains difficult. Existing solutions for CNN partitioning often do not clearly address recovery when an execution node or communication link fails.\\

Due to those constraints, it is challenging to satisfy the demanding latency requirements solely with the onboard processing unit of the device. A widespread methodology to reduce the end-to-end delay is called collaborative inference. It is also known as device edge inference. A DNN is split at a chosen layer. Then the device transfers the intermediate feature representation to an edge server, which is more capable of further calculations. It is important to select an appropriate partition point because different layer segments may vary between devices in terms of computation cost and the volume of data transmitted. These variations directly affect the communication latency\cite{Mu2024FineGrain}. \\

In real life, operational efficiency should be considered during decision making. Depending on node capabilities and power attributes, the heterogeneity of the nodes distributed execution across heterogeneous nodes can trade delay against energy usage. Furthermore, mobility plays a huge role when an inference model is equipped with the UAV-based edge infrastructures. Channel conditions and availability of nodes may change with time. Therefore, prediction-aware planning rather than static offloading can be suggested for such a system\cite{EnergyEffUAV}.\\


This thesis provides a solution to the mentioned challenges in a setting of an agricultural field monitoring scenario by using a swarm of UAVs and an edge computer located in the middle of the field. U2U and U2E communication link quality often varies with the position and speed of the UAV. Therefore, the communication cost should be calculated and updated in real-time. The CNN inference model is used to analyze an image of the field. The workload spreads across the UAV network as the sequential CNN segments may run on different nodes. Then the feature tensors will be offloaded within the boundaries of each segment. In addition to the selection of suitable execution nodes, this thesis introduces a hierarchical fallback logic. When the edge coordinator is unavailable or when UAVs fail or become unresponsive, the fallback logic triggers and continues the inference procedure.


\section{Motivation}


Focusing on intelligent crop monitoring is a meaningful research direction. It is challenging to execute a CNN model completely on a small UAV because of constrained onboard energy, limited computing resources, and variable communication quality. To address this issue, collaborative inference is a suitable solution because it allows a model to be executed across multiple nodes. However, the efficiency of this approach depends on the selected split point and the amount of intermediate data that must be transmitted \cite{second_ref}.\\

The motivation for this thesis is to improve UAV swarm and edge coordinated monitoring through adaptive CNN partitioning, improved segment placement, and structured recovery behavior. The aim is to reduce the delay and improve efficiency while addressing the gap explained in the Section \ref{Research_Gap} with detailed fallback logic. The thesis furthermore motives to improve the robustness of the algorithm by defining guardrails for deadline, connectivity, mobility, battery, responsiveness, and stability. By following these rules, the inference procedure can continue even when other nodes fail or communication links degrade.



\section{Problem statement}

The direct question initially occurring in this scenario is whether the UAVs can capture images and transfer them to the edge computer without onboard processing, allowing the full CNN to be executed at the edge. Which initially appears as the simplest idea. Even though the edge node has higher computational capacity, this approach requires transferring the complete raw image over the wireless link and can increase the communication delay. By integrating collaborative inference can avoid this by processing the first CNN layers internally and then transmitting an intermediate feature representation only at a selected split point. The feature tensors are usually smaller than the original image. Therefore, it reduces transmission time and lowers the round-trip delay\cite{HybridFastInference2023Zhang}\cite{Mu2024FineGrain}. This effect is also reflected in the simulation results discussed in Section \ref{simu_discussion}.\\



The scenarios that employ collaborative inference, Latency is the primary consideration when selecting the most suitable and efficient split point for a CNN \cite{Mu2024FineGrain}. Therefore, the challenge is not only to determine which node to offload, but also depends on a split decision in order to get the lowest possible latency. Furthermore, the hard deadline requirements, U2U and U2E link throughput, node workload, and the mobility details of UAVs should also be considered while calculating optimal node offloading and splitting decisions  \cite{Suganya2024DTOE}\cite{demir2024}. Since the system is based on energy-constrained UAVs, the scheduler cannot reduce latency without considering energy usage. Therefore, computation energy, transmission energy, and battery constraints must also be considered to keep the execution feasible.\\

Finally, the section \ref{Research_Gap} reflects that existing studies which allignes to this thesis are lacking a reliable failover and fallback mechanism. A robust failover mechanism is crucial in systems that rely on multiple distributed collaborative nodes, because there is a possibility of node failures, link degradation, etc.


\section{Objectives}
Based on the above problem statement, the objectives for addressing these issues are stated as follows.

\begin{enumerate}
	\item Design a split inference architecture that partitions a CNN into sequential segments and distributes them between the capture UAV, peer UAVs, and the edge coordinator.
	
	\item Minimize predicted end-to-end latency and reduce deadline misses under a hard deadline $\Delta t$.
	
	\item Predict the movement of the UAV and link delay over a short window to make sure the plan remains valid during execution.
	
	\item Add energy awareness and design battery guardrails for all participating UAVs.
	
	\item Ensure plan stability with re-planning with anti-thrashing controls.
	
	\item Provide continuity via re-routing and a structure failover mechanism when events of transfer failure, or when peers fail.
	
	\item Design a simulation to prove that the system can reduce latency in each scenario.
\end{enumerate}